
\section{\textbf{\texttt{write()}} ORM Method}

\begin{frame}[fragile]{\textbf{\texttt{write()}} Method}
	\begin{itemize}
		\item The \textbf{\texttt{write()}} method is an ORM method used to \textbf{update} one or more existing records in an Odoo model.
		\item Unlike \texttt{create()}, \texttt{write()} works on an existing recordset.
		\item The \texttt{write()} method returns a boolean value: \texttt{True} when the write succeeds.
\begin{block}{Method Syntax}
	\begin{itemize}
		\item The standard syntax of the write() method is:
\begin{lstlisting}
def write(self, vals):
	return super().write(vals)
\end{lstlisting}
		\item This method has three important parts:
		\begin{enumerate}
			\item \textbf{\texttt{write}}
			\item \textbf{\texttt{self}}
			\item \textbf{\texttt{vals}}
		\end{enumerate}
	\end{itemize}
\end{block}
	\end{itemize}
\end{frame}

\begin{frame}[fragile]{Breaking Down the Method Signature}
	\begin{block}{}
		\begin{itemize}
			\item \textbf{\texttt{write}}: The ORM method responsible for updating records.
			\item \textbf{\texttt{self}}: A recordset of one or more records that will be updated.
			      If we browse student ID 15, \texttt{self} represents that student record.
\begin{lstlisting}
student = self.env['student.student'].browse(15)
print(student)
\end{lstlisting}
			Output is:
\begin{lstlisting}
student.student(15,)
\end{lstlisting}
			\item If \texttt{self} contains many records, \texttt{write()} updates \textbf{all} of them with the same values.
		\end{itemize}
	\end{block}
\end{frame}

\begin{frame}[fragile]{Breaking Down the Method Signature}
	\begin{block}{}
		\begin{itemize}
			\item \textbf{\texttt{vals}}: A dictionary containing the field values to update.
\begin{lstlisting}
vals = {
	'name': 'Ahmed Muhumed',
	'email': 'ahmed@example.com',
	'age': 23,
}
\end{lstlisting}
			\begin{itemize}
				\item Each \textbf{key} is a field name on the model.
				\item Each \textbf{value} is the new data to store.
				\item Fields that are not listed in \texttt{vals} stay unchanged.
				\item The ORM maps each key to its corresponding database column.
			\end{itemize}
		\end{itemize}
	\end{block}
\end{frame}

\begin{frame}[fragile]{What Does \texttt{write()} Return?}
	\begin{block}{}
		\begin{itemize}
			\item The \textbf{\texttt{write()}} method returns \texttt{True} on success.
\begin{lstlisting}
student = self.env['student.student'].browse(15)
result = student.write({
	'name': 'Ahmed Muhumed',
})
print(result)
\end{lstlisting}
			Output is:
\begin{lstlisting}
True
\end{lstlisting}
			\item The record itself is updated in the database.
			\item You can continue using the same recordset after writing.
		\end{itemize}
	\end{block}
\end{frame}

\begin{frame}[fragile]{Updating One Record vs Many Records}
	\begin{enumerate}
		\item \textbf{One record:}
\begin{lstlisting}
student = self.env['student.student'].browse(15)
student.write({'age': 23})
\end{lstlisting}
		\item \textbf{Many records (same values for all):}
\begin{lstlisting}
students = self.env['student.student'].browse([15, 16, 17])
students.write({'active': True})
\end{lstlisting}
	\end{enumerate}
	\begin{itemize}
		\item One \texttt{write()} call on a multi-recordset is more efficient than looping and writing each record separately.
		\item Inverse and computed fields may also be recomputed after \texttt{write()}.
	\end{itemize}
\end{frame}

\begin{frame}[fragile]{Method Call and Override Pattern}
\begin{block}{Method Call}
	\begin{itemize}
		\item A method call is when we \textbf{execute} the method.
\begin{lstlisting}
student.write({
	'name': 'Ahmed Muhumed',
	'email': 'ahmed@example.com',
})
\end{lstlisting}
		\item Here, \textbf{\texttt{.write(\{...\})}} is the \textbf{method call}.
	\end{itemize}
\end{block}
\begin{block}{Common Override Pattern}
\begin{lstlisting}
def write(self, vals):
	# custom logic before saving
	if 'email' in vals:
		vals['email'] = vals['email'].lower()
	return super().write(vals)
\end{lstlisting}
\end{block}
\end{frame}
